\chapter{Введение}

\section*{Как устроена современная электроника}
\addcontentsline{toc}{section}{Как устроена современная электроника}

Большинство современных электронных устройств представляют собой компьютеры
независимо от того, насколько сложные функции они при этом выполняют.
Например, <<умная>> лампочка загорается, когда становится темно.
Лифт едет после нажатия кнопки. А мобильный телефон показывает смешную
картинку, если человек нажал на правильное место на экране.

Все эти вещи кажутся очень разными, а на самом деле они очень похожи.
Внутри каждого такого устройства есть маленький или большой компьютер.
Компьютеры отличаются от других электрических и электронных устройств тем,
что работают по программе, которую можно изменить, не переделывая всё устройство
целиком.

Программа похожа на кулинарный рецепт: она состоит из последовательности
простых шагов. Если взять другую последовательность, то получится другое блюдо,
несмотря на одинаковые кастрюли и сковородки.

Сами же простые шаги --- это операции с числами: сложение, вычитание, умножение,
сравнение и т.п.
Например, одним из чисел может быть температура воздуха (которую нужно измерить),
а другим --- напряжение на обогревателе, который нужно включить, если слишком холодно.
Программисты записывают эти вычисления и процесс принятия решений в виде программы,
а компьютеры затем программы выполняют.

Все данные, которое есть в компьютерах (телефонах, игровых приставках) представляются
в виде чисел. Например, картинка это набор точек, а каждая точка описывается одним или
несколькими числами, определяющими её цвет. В свою очередь видео --- это множество картинок, дополненных
звуковой дорожкой. Звук же можно закодировать последовательностью чисел, описывающих
уровень напряжения на воспроизводящем динамике.

Отличаются компьютеры между собой тем, сколько чисел (то есть данных) они могут запоминать
и с какой скоростью обрабатывать. Ну и конечно, такими незначительными деталями,
как количество пикселей на экране или число кнопок у <<мыши>>.


\section*{Из чего состоит компьютер}
\addcontentsline{toc}{section}{Из чего состоит компьютер}

% Внутри компьютеров спрятано много маленьких вычислителей, работающих одновременно
% или соединённых в цепочки. Но как компьютеру удаётся с помощью них делать что-то полезное?

Главное, что отличает компьютеры от других механизмов (или электронных устройств вроде часов) ---
это программы. В компьютере за работу по программе отвечает процессор.
Именно он последовательно её <<читает>> и делает всё, что там написано.
Фрагменты программы (шаги, инструкции, команды) обычно очень простые. Например, сложить два числа
или перемножить их. Потом уже из этих команд складываются программы, делающие практически что угодно
(если программа достаточно сложная).

Кроме процессора в компьютере есть ещё память. Там хранятся все данные пользователя
(музыка, фотографии, электронные письма), а также программы для обработки этих данных.
Программа обычно хранится в памяти как последовательность команд. Чтобы двигаться по этим командам по мере их выполнения,
нужно запоминать номер текущей. За это отвечает счётчик инструкций (program counter, instruction pointer).
Счётчик инструкций обычно находится внутри процессора.

Процессор <<смотрит>> в память на текущую инструкцию и активирует один из своих модулей, отвечающий за её выполнение.
Это может быть арифметическая операция, переход к другой части программы (тогда в счётчик инструкций попадает новое значение),
передача данных на внешнее устройство (вывод пикселя на экран). После этого счётчик инструкций переключается на следующую
команду.

Чтобы проделывать все эти шаги поочерёдно, используется тактовый генератор. Он с заданной частотой посылает
сигналы на все остальные модули, в результате чего они активируются один за другим. Например, сначала
процессор смотрит в память, потом декодирует инструкцию, то есть, разбирается с тем,
какие именно действия нужно будет выполнять дальше, а после выполнения этих действий
увеличивает счётчик инструкций.

Все числа (данные), с которыми работает процессор хранятся в памяти. Он может считывать их оттуда и записывать новые.
Но внешняя (по отношению к процессору) память обычно работает не слишком быстро, поэтому в процессоре есть несколько
регистров, используемых для вычислений. Например, одна команда процессора может сложить числа из двух регистров и записать результат
в третий. А потом другая команда отправит значение этого регистра во внешнюю память.

Мы будем строить собственный компьютер с помощью релейного конструктора.
То есть формально это будет компьютер, но он будет настолько простой, что можно даже
называть его процессором. А в качестве промежуточных этапов у нас будут получаться
ещё более простые вычислители.

В целом, компьютер можно сделать из чего угодно: из шестерёнок, реле, транзисторов,
микросхем.
Электромагнитные реле хороши тем, что сразу понятно, что они делают.
Современные компьютеры построены на транзисторах, которые упакованы в микросхемы,
но с ними не так интересно играть. Например, реле можно переключить простой отвёрткой,
а потом оценить, что изменилось.

Самое главное, что современные процессоры и компьютеры концептуально
почти не отличаются от содержимого наших занятий.
Мы будем создавать следующие части компьютера (и процессора):
\begin{enumerate}
    \item Регистры для хранения данных.
    \item Вычислительные модули для реализации команд.
    \item Память программ.
    \item Счётчик инструкций.
    \item Логические схемы для декодирования команд и выбора правильных исполнительных модулей.
    \item Тактовый генератор.
\end{enumerate}

Такой компьютер сможет выполнять простые программы, например, найти остаток от деления одного числа на другое,
используя только сложение и вычитание. А если в какой-то момент реле станет недостаточно,
полученные знания можно будет применить для разработки более сложного процессора
на современных технологиях или при программировании уже готовых устройств.

\section*{Для чего эта книга и как ей пользоваться}
\addcontentsline{toc}{section}{Для чего эта книга и как ей пользоваться}
